% \chapter{Phương pháp chèn bảng}
% \label{chap:chap3-table}

% Chương này trình bày một số kỹ thuật chèn bảng với các cột, dòng merge nhiều ô. Các cách format bảng cơ bản không trình bày. Trong quá trình sử dụng, có thể cập nhật các chỉ dẫn rõ ràng cho từng câu lệnh nghĩa là gì trong một bảng.

% \section{Chèn bảng có nhiều hàng, cột}

% Bảng \ref{tab:student_statistics} mô tả một bảng mẫu bao gồm đầy đủ các yếu tố như cột nhiều hàng, hàng nhiều cột. Các kỹ thuật sử dụng trong bảng này cơ bản đã bao gồm hết các thao tác thông thường của việc tạo bảng.

% This is an example document to demonstrate the use of glossaries. 
% We mention \gls{ssl} and \gls{api} here.

% \begin{table}[ht]
%   \centering
%   \caption{Comparison of statistics among students in different faculties}
%   \label{tab:student_statistics}
%   \begin{tabular}{|C{.2\textwidth} | C{.2\textwidth} | C{.6\textwidth}|}
%     \hline
%     \multicolumn{2}{|c|}{\textbf{Faculty}} & \multirow{2}{*}{\textbf{Attributes}} \\
%     \cline{1-2}
%     \textbf{Faculty A} & \textbf{Faculty B} & \\
%     \hline
%     \multirow{3}{*}{Physics} & Average Score & Students in Faculty A have an average score of 85, while those in Faculty B have an average score of 78. \\
%     \cline{2-3}
%     & Physical Education & Faculty A emphasizes physical education with mandatory sports classes, resulting in higher fitness levels compared to Faculty B. \\
%     \cline{2-3}
%     & Book Reading & Students in Faculty A are encouraged to read books related to physics, which has positively impacted their understanding of the subject. \\
%     \hline
%     \multirow{3}{*}{\shortstack[C]{Computer \\ Science}} & Average Score & Faculty A's students have an average score of 90, whereas Faculty B's students have an average score of 82. \\
%     \cline{2-3}
%     & Culture Activity & Faculty A organizes regular cultural activities, enhancing students' cultural awareness and social skills. \\
%     \cline{2-3}
%     & Internship Opportunities & Faculty A provides more internship opportunities, leading to better practical knowledge among students. \\
%     \hline
%   \end{tabular}
% \end{table}

% \section{Chèn bảng có nhiều trang}

% Trong một số trường hợp cần chèn bảng với quá nhiều dòng, số liệu của chúng sang trang khác. Ví dụ như trường hợp ta cần mô tả bảng số liệu của quá trình training qua từng vòng, thì sử dụng \textit{longtable} là sự lựa chọn thông minh. Bảng \ref{tab:chap3-longtable-example} mô tả một bảng như vậy.

% \begin{longtable}[htbp]{|c|| c|| c|| c|| c|| c|| c|} 
%   \caption{Mô tả bảng dài với nhiều số liệu}
%   \label{tab:chap3-longtable-example} \\
%   \hline
%   \hline
%   {Data} & {$\beta_{UP}$} & {$ \beta_{LOW}$} & {$ \beta_{TEL}$} & {$R_{UP}$} & {$R_{LOW}$}  & {$R_{TEL}$} \\
%   \hline
%   \endfirsthead
%   \hline
%   {Data} & {$\beta_{UP}$} & {$ \beta_{LOW}$} & {$ \beta_{TEL}$} & {$R_{UP}$} & {$R_{LOW}$}  & {$R_{TEL}$} \\
%   \hline
%   \endhead
%   \hline \multicolumn{7}{|r|}{{Continued on next page}} \\ \hline
%   \endfoot
%   \hline
%   \multicolumn{7}{|c|}{{End of Table}} \\ \hline
%   \endlastfoot

%   200604 & 0.0437 & 0.0087 & -0.0963 & 0.17 & 0.05 & -0.46 \\
%   200605 & 0.0317 & 0.0478 & -0.1136 & 0.20 & 0.40 & -0.61 \\
%   200606 & -0.0699 & -0.0418 & -0.0420 & -0.21 & -0.16 & -0.005 \\
%   200607 & -0.0783 & -0.0483 & -0.0868 & -0.24 & -0.19 & -0.25 \\
%   200608 & -0.0551 & -0.0734 & -0.1932 & -0.27 & -0.41 & -0.63 \\
%   201410 & -0.0609 & -0.0849 & -0.3131 & -0.06 & -0.06 & -0.38 \\
%   201411 & -0.0279 & -0.0849 & -0.2353 & -0.03 & -0.09 & -0.29 \\
%   201412 & 0.0884 & -0.0252 & -0.1411 & 0.08 & -0.02 & -0.12 \\ 
%   200604 & 0.0437 & 0.0087 & -0.0963 & 0.17 & 0.05 & -0.46 \\
%   200605 & 0.0317 & 0.0478 & -0.1136 & 0.20 & 0.40 & -0.61 \\
%   200606 & -0.0699 & -0.0418 & -0.0420 & -0.21 & -0.16 & -0.005 \\
%   200607 & -0.0783 & -0.0483 & -0.0868 & -0.24 & -0.19 & -0.25 \\
%   200608 & -0.0551 & -0.0734 & -0.1932 & -0.27 & -0.41 & -0.63 \\
%   201410 & -0.0609 & -0.0849 & -0.3131 & -0.06 & -0.06 & -0.38 \\
%   201411 & -0.0279 & -0.0849 & -0.2353 & -0.03 & -0.09 & -0.29 \\
%   201412 & 0.0884 & -0.0252 & -0.1411 & 0.08 & -0.02 & -0.12 \\ 
%   200604 & 0.0437 & 0.0087 & -0.0963 & 0.17 & 0.05 & -0.46 \\
%   200605 & 0.0317 & 0.0478 & -0.1136 & 0.20 & 0.40 & -0.61 \\
%   200606 & -0.0699 & -0.0418 & -0.0420 & -0.21 & -0.16 & -0.005 \\
%   200607 & -0.0783 & -0.0483 & -0.0868 & -0.24 & -0.19 & -0.25 \\
%   200608 & -0.0551 & -0.0734 & -0.1932 & -0.27 & -0.41 & -0.63 \\
%   201410 & -0.0609 & -0.0849 & -0.3131 & -0.06 & -0.06 & -0.38 \\
%   201411 & -0.0279 & -0.0849 & -0.2353 & -0.03 & -0.09 & -0.29 \\
%   201412 & 0.0884 & -0.0252 & -0.1411 & 0.08 & -0.02 & -0.12 \\ 
%   200604 & 0.0437 & 0.0087 & -0.0963 & 0.17 & 0.05 & -0.46 \\
%   200605 & 0.0317 & 0.0478 & -0.1136 & 0.20 & 0.40 & -0.61 \\
%   200606 & -0.0699 & -0.0418 & -0.0420 & -0.21 & -0.16 & -0.005 \\
%   200607 & -0.0783 & -0.0483 & -0.0868 & -0.24 & -0.19 & -0.25 \\
%   200608 & -0.0551 & -0.0734 & -0.1932 & -0.27 & -0.41 & -0.63 \\
%   201410 & -0.0609 & -0.0849 & -0.3131 & -0.06 & -0.06 & -0.38 \\
%   201411 & -0.0279 & -0.0849 & -0.2353 & -0.03 & -0.09 & -0.29 \\
%   201412 & 0.0884 & -0.0252 & -0.1411 & 0.08 & -0.02 & -0.12 \\ 
%   200604 & 0.0437 & 0.0087 & -0.0963 & 0.17 & 0.05 & -0.46 \\
%   200605 & 0.0317 & 0.0478 & -0.1136 & 0.20 & 0.40 & -0.61 \\
%   200606 & -0.0699 & -0.0418 & -0.0420 & -0.21 & -0.16 & -0.005 \\
% \end{longtable}

% Ngoài ra, trong một vài trường hợp có thể sử dụng bảng ngang. Nhưng xét thấy ít khi sử dụng vậy nên tác giả không thêm vào.

% % \begin{sidewaystable}
% %   \centering
% %   \caption{Comparison}
% %   \begin{tabular}{l*{6}{c}r}
% %   \hline
% %   Names  & A & B & C & D \\
% %   \hline
% %   \hline
% %   Jobs   & A & B & C & D \\
% %   \hline
% %   Types   & A & B & C & D \\
% %   \hline
% %   \end{tabular}    
  
% %   \end{sidewaystable}

\chapter{Kết quả đạt được và hạn chế}
\label{chap:chap3-result-limit}

Chương này trình bày các kết quả đạt được của đề tài sau quá trình thiết kế, triển
khai và kiểm thử hệ thống, đồng thời chỉ ra những hạn chế còn tồn tại. Các nội dung
được tổng hợp dựa trên toàn bộ quá trình hiện thực hoá hệ thống, triển khai hạ tầng,
cấu hình các công cụ vận hành và kết quả kiểm thử đã được trình bày trong Chương~2.

\section{Kết quả đạt được}
\label{sec:chap3-results}

Sau quá trình triển khai và đánh giá, đề tài đã đạt được các kết quả chính như sau:

\subsection{Hiện thực thành công hệ thống Kubernetes cho ứng dụng microservices}

Đề tài đã thiết kế và triển khai thành công một hệ thống Kubernetes cho ứng dụng
microservices trên nền tảng điện toán đám mây Amazon Web Services. Hệ thống được
triển khai theo kiến trúc cloud-native, tuân thủ các nguyên tắc của hệ thống phân
tán hiện đại, bao gồm phân tách môi trường, khả năng mở rộng và khả năng chịu lỗi.

Ứng dụng microservices được triển khai với đầy đủ các thành phần quản lý cấu hình,
khám phá dịch vụ và điều phối truy cập, cho phép các dịch vụ hoạt động độc lập nhưng
vẫn phối hợp chặt chẽ trong quá trình xử lý request. Việc triển khai ứng dụng trên
Kubernetes giúp hệ thống có khả năng tự động quản lý vòng đời Pod, phân bố workload
và khôi phục khi xảy ra sự cố.

\subsection{Triển khai đa môi trường với kiến trúc nhất quán}

Hệ thống được triển khai trên hai môi trường độc lập gồm môi trường kiểm thử
(\textit{Staging}) và môi trường sản xuất mô phỏng (\textit{Production}). Mỗi môi
trường được thiết kế với mức độ phức tạp và mục tiêu khác nhau nhưng vẫn đảm bảo
tính nhất quán về kiến trúc tổng thể.

Môi trường staging sử dụng Kubernetes k0s để mô phỏng hành vi vận hành của Kubernetes
trong điều kiện thực tế, phục vụ cho việc kiểm thử cấu hình và đánh giá kiến trúc.
Môi trường production sử dụng Amazon EKS nhằm tận dụng các dịch vụ managed của AWS,
đảm bảo độ ổn định và tính sẵn sàng cao cho hệ thống.

Cách tiếp cận này giúp quá trình phát triển, kiểm thử và đánh giá hệ thống được thực
hiện có kiểm soát, đồng thời phản ánh mô hình triển khai phổ biến trong các hệ thống
doanh nghiệp.

\subsection{Tự động hoá triển khai hạ tầng và cấu hình hệ thống}

Toàn bộ hạ tầng của hệ thống được định nghĩa và triển khai bằng Terraform theo mô
hình \textit{Infrastructure as Code}. Việc sử dụng Terraform cho phép tự động hoá
quá trình khởi tạo tài nguyên, đảm bảo tính nhất quán giữa các môi trường và giảm
thiểu sai sót do thao tác thủ công.

Bên cạnh đó, các thành phần cấu hình hệ thống như OpenVPN, Prometheus, Grafana,
Loki và Promtail được triển khai tự động bằng Ansible. Việc kết hợp Terraform và
Ansible giúp hệ thống có khả năng tái triển khai nhanh chóng, dễ mở rộng và phù
hợp với yêu cầu vận hành trong môi trường cloud-native.

\subsection{Xây dựng hệ thống quan sát cơ bản cho Kubernetes}

Đề tài đã xây dựng thành công một hệ thống quan sát (observability) ở mức cơ bản,
bao gồm giám sát metric và ghi log tập trung. Prometheus được sử dụng để thu thập
metric hạ tầng và Kubernetes, trong khi Grafana đóng vai trò trực quan hoá dữ liệu
và cung cấp dashboard tổng quan cho hệ thống.

Hệ thống logging được xây dựng dựa trên Loki và Promtail, cho phép thu thập log từ
toàn bộ Pod và container trong cụm Kubernetes về một điểm tập trung. Việc kết hợp
giữa giám sát metric và ghi log giúp người vận hành có cái nhìn tổng thể về trạng
thái hệ thống, hỗ trợ quá trình phân tích và đánh giá hành vi vận hành.

\subsection{Đánh giá khả năng mở rộng và tự phục hồi của hệ thống}

Thông qua các kịch bản kiểm thử tải và mô phỏng sự cố, đề tài đã đánh giá được khả
năng mở rộng và tự phục hồi của hệ thống. Kết quả kiểm thử cho thấy hệ thống có khả
năng xử lý tải người dùng đồng thời lớn và tự động mở rộng số lượng Pod thông qua
Horizontal Pod Autoscaler.

Trong kịch bản mô phỏng mất node, Kubernetes và Amazon EKS đã thể hiện khả năng tự
phục hồi hiệu quả. Các Pod bị mất được tự động khởi tạo lại và phân bố sang các node
còn lại, đồng thời node mới được tạo ra để tái cân bằng tải. Trong suốt quá trình
này, hệ thống vẫn duy trì dịch vụ ổn định và không ghi nhận gián đoạn nghiêm trọng
từ phía người dùng.

\subsection{Khả năng mô phỏng và đánh giá hệ thống phân tán trong điều kiện thực tế}

Thông qua việc triển khai hệ thống trên hạ tầng cloud và thực hiện các kịch bản
kiểm thử có kiểm soát, đề tài đã mô phỏng được các tình huống vận hành phổ biến
trong hệ thống phân tán thực tế như tăng tải người dùng, mất node và quá trình
phục hồi tự động.

Các kết quả thu thập từ hệ thống giám sát và công cụ kiểm thử cho phép đánh giá
trực quan hành vi của hệ thống ở cả mức hạ tầng và mức ứng dụng. Điều này giúp
làm rõ vai trò của Kubernetes trong việc quản lý tài nguyên, phân phối workload
và duy trì tính sẵn sàng của dịch vụ trong môi trường cloud-native.

Việc mô phỏng và đánh giá hệ thống trong các kịch bản khác nhau không chỉ giúp
kiểm chứng tính đúng đắn của kiến trúc đề xuất, mà còn cung cấp cơ sở thực tiễn
cho việc phân tích, so sánh và cải tiến hệ thống trong các nghiên cứu hoặc triển
khai tiếp theo.


\section{Những hạn chế cần cải thiện trong tương lai}
\label{sec:chap3-limitations}

Bên cạnh các kết quả đạt được, đề tài vẫn còn tồn tại một số hạn chế nhất định do
giới hạn về thời gian, kiến thức và phạm vi thực hiện.

\subsection{Hạn chế về mức độ hoàn thiện của hệ thống quan sát}

Hệ thống observability trong đề tài mới dừng lại ở mức cơ bản, tập trung vào giám
sát metric và ghi log. Các thành phần nâng cao như distributed tracing, phân tích
log nâng cao hoặc cảnh báo thông minh chưa được triển khai đầy đủ.

Ngoài ra, hệ thống chưa xây dựng các cơ chế cảnh báo tự động (alerting) gắn với
các chỉ số SLO ( Service-Level Objective) hoặc SLA ( Service-Level Agreement), do đó khả năng phát hiện và phản ứng sớm với sự cố vẫn còn
hạn chế.

\subsection{Hạn chế về bảo mật và công khai dịch vụ}

Ứng dụng trong đề tài chưa được công khai hoàn toàn ra Internet với đầy đủ các cơ
chế bảo mật nâng cao. Việc triển khai HTTPS, quản lý chứng chỉ TLS tự động, Web
Application Firewall (WAF) và các chính sách bảo mật lớp ứng dụng chưa được hiện
thực trong phạm vi đề tài.

Hệ thống hiện tại chủ yếu phục vụ cho mục đích mô phỏng, kiểm thử và đánh giá kiến
trúc, chưa hướng đến vận hành trong môi trường sản xuất thực tế với yêu cầu bảo
mật nghiêm ngặt.

\subsection{Hạn chế về phạm vi tự động hoá và tối ưu vận hành}

Mặc dù hạ tầng và cấu hình hệ thống đã được tự động hoá ở mức cơ bản, các cơ chế
CI/CD, GitOps hoặc tự động tối ưu tài nguyên chưa được triển khai. Việc cập nhật
ứng dụng và cấu hình hệ thống vẫn cần sự can thiệp thủ công trong một số bước.

Ngoài ra, hệ thống chưa thực hiện các bài kiểm thử hiệu năng chuyên sâu hoặc đánh
giá chi phí vận hành dài hạn, do đó chưa phản ánh đầy đủ các yếu tố tối ưu trong
môi trường production thực tế.

Nhìn chung, các hạn chế nêu trên xuất phát từ phạm vi và mục tiêu của đề tài,
trong đó trọng tâm là hiện thực và đánh giá kiến trúc hệ thống phân tán trên nền
tảng Kubernetes. Những nội dung chưa được triển khai đầy đủ chủ yếu liên quan đến
các thành phần nâng cao và yêu cầu vận hành trong môi trường production thực tế.

Các hạn chế này không làm ảnh hưởng đến mục tiêu chính của đề tài, nhưng đóng vai
trò quan trọng trong việc xác định phạm vi và định hướng nghiên cứu, phát triển
ở các giai đoạn tiếp theo.
